\documentclass[preprint,12pt]{elsarticle}
\usepackage{amsmath,amssymb}
\usepackage{graphicx}
\usepackage{booktabs}
\usepackage{multirow}
\usepackage{tabularx}
\usepackage{array}
\usepackage{booktabs}
\usepackage{float}
\usepackage{enumitem}
\usepackage{booktabs}
\usepackage{makecell}
\usepackage{ragged2e}
\usepackage{placeins}
\usepackage{url}
\usepackage{microtype}
\usepackage{subcaption}
\usepackage{enumitem}
\usepackage{siunitx}
\graphicspath{{figures/}}

\journal{Journal manuscript draft - target journal to be selected}

\begin{document}

\begin{frontmatter}

\title{Relative Amplitude TD Exploration: Reward-Scale Robust Adaptive Exploration and Approximator-Dependent Performance in Value-Based Reinforcement Learning}

\author[bracu]{Ihzaz Abrar Sadik}
%\author{Co-authors to be inserted}
\affiliation[bracu]{organization={Department of Computer Science and Engineering, BRAC University}, city={Dhaka}, country={Bangladesh}}

\begin{abstract}

\end{abstract}

\begin{keyword}
reinforcement learning \sep exploration \sep temporal-difference error \sep function approximation \sep VDBE \sep Double DQN \sep reward scaling \sep robust evaluation
\end{keyword}

\end{frontmatter}

\section{Introduction}



\section{Background and Related Work}
\subsection{Value-based control and TD error}


\subsection{Exploration schedules and adaptive exploration}


\subsection{Evaluation under seed variability}

\section{Relative Amplitude TD Exploration}
\subsection{Running amplitudes}

\subsection{Positive reward-scale property}


\subsection{Computational cost and neural integration}


\section{Experimental Methodology}
\subsection{Environments and interaction budgets}
Four environments with continuous state spaces, discrete actions but varying structure have been used. They all give the exploration rules different kinds of feedback. This helps in assessing the exploration methods across a diverse range of scenarios. The `Gymnasium' package in Python, contains these environments and were used by parsing them from it. Summary of all the environments are listed in table \ref{tab:environments}. 

\begin{table}[H]
    \centering
    \caption{Summary of the benchmark environments.}
    \label{tab:environments}
    \renewcommand{\arraystretch}{1.15}
    \setlength{\tabcolsep}{7pt}

    \begin{tabular}{lcccc}
        \toprule
        \textbf{Environment} &
        \textbf{Obs. Dim.} &
        \textbf{State Dime.} &
        \textbf{$|A|$} &
        \textbf{Final Budget} \\
        \midrule

        MountainCar-v0 & 2 & 2 & 3 & 150,000 \\
        CartPole-v1    & 4 & 4 & 2 & 100,000 \\
        Acrobot-v1     & 6 & 4 & 3 & 100,000 \\
        LunarLander-v3 & 8 & 8 & 4 & 300,000 \\

        \bottomrule
    \end{tabular}
\end{table}

\noindent\textbf{Environment Details}

\begin{enumerate}[leftmargin=*, itemsep=0.8em]

    \item
    \begin{minipage}[t]{\linewidth}
    \textbf{MountainCar-v0}
    \begin{itemize}[leftmargin=1.5em, itemsep=0.2em, topsep=0.2em]
        \item \textbf{Brief Intro:} An underpowered car in a valley between two hills, must reach a flag on the right hill. The engine is weaker than gravity.
        \item \textbf{Reward Structure:} -1 per step until goal
        \item \textbf{Significance:} Success is sparse, exploration needs to discover momentum in the opposite direction to the goal
    \end{itemize}
    \end{minipage}

    \item
    \begin{minipage}[t]{\linewidth}
    \textbf{CartPole-v1}
    \begin{itemize}[leftmargin=1.5em, itemsep=0.2em, topsep=0.2em]
        \item \textbf{Brief Intro:} A pole is hinged on a cart; it can be pushed right or left to keep the pole upright.
        \item \textbf{Reward Structure:} +1 for each surviving step
        \item \textbf{Significance:} Dense reward from the very first step
    \end{itemize}
    \end{minipage}

    \item
    \begin{minipage}[t]{\linewidth}
    \textbf{Acrobot-v1}
    \begin{itemize}[leftmargin=1.5em, itemsep=0.2em, topsep=0.2em]
        \item \textbf{Brief Intro:} A two-link pendulum hanging form a pivot, actuated only at the middle joint, must swing its tip above a given height.
        \item \textbf{Reward Structure:} -1 per step until goal
        \item \textbf{Significance:} Structured long-horizon control
    \end{itemize}
    \end{minipage}

    \item
    \begin{minipage}[t]{\linewidth}
    \textbf{LunarLander-v3}
    \begin{itemize}[leftmargin=1.5em, itemsep=0.2em, topsep=0.2em]
        \item \textbf{Brief Intro:} A lander must set down softly on a pad between two flags.
        \item \textbf{Reward Structure:} shaped+terminal
        \item \textbf{Significance:} Stochastic, high-dimensional, complex reward structure
    \end{itemize}
    \end{minipage}

\end{enumerate}

\subsection{Linear approximator}


\subsection{Double DQN}


\subsection{Exploration baselines and frozen tuning}


\subsection{Greedy evaluation}


\subsection{Robust statistics}


\section{Diagnostic Experiments}
\subsection{Experiment A: Is raw TD-error magnitude a progress signal?}


\subsection{Experiment B: VDBE's TD-independent large-$\sigma$ limit}


\subsection{VDBE fidelity ablation}

\subsection{Experiment C: positive reward rescaling}


\section{Final 30-Seed Benchmark}
\subsection{Linear Sarsa($\lambda$)}

\subsection{Double DQN}


\subsection{Approximator-dependent ranking}


\section{Discussion}
\subsection{Scale robustness is not performance optimality}

\subsection{Why might the approximator matter?}


\subsection{What the VDBE diagnostics add}


\subsection{MountainCar under Double DQN}

\section{Limitations}


\section{Conclusion}


\section*{Data and Reproducibility Statement}

\begin{thebibliography}{99}
\bibitem{sutton2018} R. S. Sutton and A. G. Barto, \emph{Reinforcement Learning: An Introduction}, 2nd ed. Cambridge, MA: MIT Press, 2018.
\bibitem{tokic2010} M. Tokic, ``Adaptive $\epsilon$-greedy exploration in reinforcement learning based on value differences,'' in \emph{KI 2010: Advances in Artificial Intelligence}, R. Dillmann, J. Beyerer, U. Hanebeck, and T. Schultz, Eds., Lecture Notes in Computer Science, vol. 6359. Berlin, Heidelberg: Springer, 2010, pp. 203--210. doi: 10.1007/978-3-642-16111-7\_23.
\bibitem{auer2002} P. Auer, N. Cesa-Bianchi, and P. Fischer, ``Finite-time analysis of the multiarmed bandit problem,'' \emph{Machine Learning}, vol. 47, pp. 235--256, 2002.
\bibitem{mnih2015} V. Mnih, K. Kavukcuoglu, D. Silver, A. A. Rusu, J. Veness, M. G. Bellemare, A. Graves, M. Riedmiller, A. K. Fidjeland, G. Ostrovski, S. Petersen, C. Beattie, A. Sadik, I. Antonoglou, H. King, D. Kumaran, D. Wierstra, S. Legg, and D. Hassabis, ``Human-level control through deep reinforcement learning,'' \emph{Nature}, vol. 518, pp. 529--533, 2015. doi: 10.1038/nature14236.
\bibitem{hasselt2016} H. van Hasselt, A. Guez, and D. Silver, ``Deep reinforcement learning with Double Q-learning,'' in \emph{Proceedings of the AAAI Conference on Artificial Intelligence}, vol. 30, no. 1, 2016, pp. 2094--2100. doi: 10.1609/aaai.v30i1.10295.
\bibitem{bellemare2016} M. G. Bellemare, S. Srinivasan, G. Ostrovski, T. Schaul, D. Saxton, and R. Munos, ``Unifying count-based exploration and intrinsic motivation,'' in \emph{Advances in Neural Information Processing Systems}, vol. 29, 2016, pp. 1471--1479.
\bibitem{osband2016} I. Osband, C. Blundell, A. Pritzel, and B. Van Roy, ``Deep exploration via bootstrapped DQN,'' in \emph{Advances in Neural Information Processing Systems}, vol. 29, 2016, pp. 4026--4034.
\bibitem{pathak2017} D. Pathak, P. Agrawal, A. A. Efros, and T. Darrell, ``Curiosity-driven exploration by self-supervised prediction,'' in \emph{Proceedings of the 34th International Conference on Machine Learning}, vol. 70, 2017, pp. 2778--2787.
\bibitem{agarwal2021} R. Agarwal, M. Schwarzer, P. S. Castro, A. C. Courville, and M. G. Bellemare, ``Deep reinforcement learning at the edge of the statistical precipice,'' in \emph{Advances in Neural Information Processing Systems}, vol. 34, 2021, pp. 29304--29320.
\end{thebibliography}

\end{document}
